% ============================================================
\section{Computational Methods}
\label{sec:methods}
% ============================================================

\subsection{Pipeline Overview}

The pipeline executes sequentially in 18 numbered steps.
Each step caches its result to disk (JSON or NumPy \texttt{.npz});
on re-run with identical inputs only missing steps are recomputed.
The high-level sequence is:

\begin{enumerate}[nosep]
  \item Download and mosaic Copernicus DEM tiles.
  \item Compute terrain derivatives (slope, aspect, horizon profile,
        TWI — cached).
  \item Download ESA WorldCover and OSM vector features.
  \item Compute \emph{preliminary} $\rzz$ and $\zzz$ at default
        $\thetav = 0.20$ and $\lw = 0$.
  \item Download SoilGrids soil properties for the footprint.
  \item Compute actual $\lw$ from clay fraction (Eq.~\eqref{eq:lw_clay}),
        recompute $\zzz$ and $\rzz$.
  \item Compute $\kT$ (3-D ray-casting over DEM).
  \item Compute $\kM$ (horizon-based muon integral).
  \item Compute $\kL$ (WorldCover + OSM, $W(r)$-weighted).
  \item Compute site neutron flux and $N_0$
        (Eqs.~\eqref{eq:Nneut_site}--\eqref{eq:N0}).
  \item Download ERA5 climatology and compute
        $\rho_{WV}$, $f_{WV}$, SWE (monthly).
  \item Download MODIS LAI and compute AGBH.
  \item Query Macrostrat for geological characterisation.
  \item Compute Topographic Wetness Index (TWI), loaded from cache
        if available.
  \item Compute optimal soil sampling plan and $\rzz$ dynamic range.
  \item Assemble the full report (PDF + JSON).
\end{enumerate}

\subsection{Terrain Processing}
\label{sec:methods:terrain}

\paragraph{Slope and aspect.}
Slope $\beta$ and aspect are derived from the DEM by the standard
central-difference gradient:
\begin{equation}
  \nabla z = \left(\frac{\partial z}{\partial x},\,
                   \frac{\partial z}{\partial y}\right),
  \qquad
  \beta = \arctan|\nabla z|.
\end{equation}
A minimum slope floor of $\SI{0.01}{\degree}$ is applied to avoid
division by zero in subsequent hydrological calculations.

\paragraph{Horizon profile.}
For each of $N_\phi = 72$ azimuth directions (5° spacing), the
maximum elevation angle $\psi(\phi)$ is computed by scanning along
the DEM at that azimuth out to $3\,\rzz$ from the site.
The result is a 72-element array $\psi(\phi_k)$ that enters both
the $\kM$ integral and the ray-casting algorithm.

\paragraph{Topographic Wetness Index.}
TWI is computed as \citep{beven1979}:
\begin{equation}
  \mathrm{TWI} = \ln\frac{A_s}{\tan\beta},
  \label{eq:twi}
\end{equation}
where $A_s$ is the specific upslope drainage area per unit contour
length, derived from the $D_\infty$ flow accumulation of the
pit-filled DEM \citep{wang2006}.
Because TWI computation is the most expensive terrain step
($\mathcal{O}(N_\mathrm{pix}\log N_\mathrm{pix})$ for large DEMs),
results are cached to disk as a \texttt{.npz} file.
The cache key is a hash of the DEM shape, pixel spacing, $\rzz$,
and summary statistics (mean and standard deviation of elevation);
if the hash matches the stored file the cached result is returned
directly.

\subsection{Topographic Correction \texorpdfstring{$\kT$}{kappa\_topo}}
\label{sec:methods:kT}

The 3-D ray-casting implementation samples $N_\theta = 18$ zenith
angles (uniformly in $\cos\theta$, $0^\circ$–$89^\circ$) and
$N_\phi = 72$ azimuth angles.
For each ray $(\theta_k, \phi_j)$:
\begin{enumerate}[nosep]
  \item March along the ray in 3-D until it intersects the DEM
        surface, recording the air path length $L_\mathrm{air}$.
  \item At the intersection, read the surface slope $\alpha$ and
        compute $L_\mathrm{soil} = \zzz / \cos\alpha$.
  \item Evaluate Eq.~\eqref{eq:ray}.
\end{enumerate}
The flat-reference integral uses a perfectly horizontal DEM at
the sensor elevation.
Both integrals are computed with the same angular sampling;
$\kT$ is their ratio.

Rays that exit the DEM bounding box without hitting terrain
(i.e., the site is on a ridge or cliff edge) contribute zero to the
numerator but retain their full flat-reference contribution in the
denominator, correctly reducing $\kT$ below 1.

\subsection{LULC Correction and Supplementary Hydrogen}
\label{sec:methods:kL}

\paragraph{WorldCover raster.}
Each 10-m WorldCover pixel within $2\,\rzz$ of the site is assigned
the hydrogen equivalence factor $f_H^{(c)}$ from Table~\ref{tab:fH}.
The radial weight $W(r_i)$ (Eq.~\eqref{eq:Wr}) is evaluated at the
pixel centroid distance $r_i$.
The correction factor $\kL$ follows Eq.~\eqref{eq:kappa_lulc}.

\paragraph{OSM vector features.}
OSM polygons (buildings, water bodies, sealed roads) are rasterised
onto the DEM grid at 1-m resolution and then downsampled to the
DEM pixel size.
For pixels with OSM overlap the WorldCover $f_H$ is replaced by the
OSM-specific value (e.g.\ $f_H = 4.0$ for permanent water,
$f_H = 0.05$ for sealed surfaces).
The uncovered area inside the $2\,\rzz$ circle contributes
\begin{equation}
  \mathrm{base\_contrib} =
    f_H^\mathrm{base} \times
    \max\!\bigl(0,\;
      \textstyle\sum_i W(r_i)\,A_i - \sum_{j\in\mathrm{OSM}} W(r_j)\,A_j
    \bigr),
  \label{eq:osm_baseline}
\end{equation}
where $f_H^\mathrm{base}$ is the WorldCover bare-soil baseline
($f_H = 0.03$) and the difference of $W$-weighted integrals
avoids the spatially averaged approximation used in simpler
implementations.

\paragraph{Reference moisture.}
All $f_H$ values and the WorldCover baseline are computed at the
user-specified initial soil moisture $\thetav^\mathrm{ref}$
(\texttt{THETA\_V\_INIT} in the configuration file, default
$\SI{0.20}{\cubic\metre\per\cubic\metre}$).
This makes $\kL$ physically consistent with the $\rzz$ and
$\zzz$ computed at the same moisture level.

\subsection{Supplementary Correction Factors}
\label{sec:methods:suppl}

These factors are not included in $\kTot$ but are reported separately
as operational correction guidance.

\paragraph{Atmospheric water vapour.}
\label{sec:methods:wv}
An increase in above-sensor water vapour reduces the epithermal
neutron flux reaching the ground.
Following \citet{rosolem2013}, the correction factor is:
\begin{equation}
  f_{WV} = \exp\!\bigl(-\alpha_{WV}\,\rho_{WV}\bigr),
  \label{eq:fWV}
\end{equation}
with $\alpha_{WV} = \SI{0.0054}{\metre\cubed\per\gram}$
\citep{zreda2012}.
Monthly climatological values of $\rho_{WV}$
(Eqs.~\eqref{eq:esat}–\eqref{eq:rhoWV}) and the corresponding
$f_{WV}$ are tabulated in the report.
Note that $f_{WV}$ is an operational (epoch-by-epoch) correction;
the climatology here serves as site-characterisation guidance.

\paragraph{Above-ground biomass hydrogen (AGBH).}
\label{sec:methods:agbh}
Vegetation intercepts and stores water; the associated hydrogen
moderates neutrons \citep{baatz2015, schron2017}.
An effective soil-moisture equivalent is estimated as:
\begin{equation}
  \mathrm{AGBH} = \mathrm{LAI} \times \SI{150}{\gram\per\metre\squared}
                  + \theta_\mathrm{litter}
  \quad[\si{\milli\metre}],
  \label{eq:agbh}
\end{equation}
where the 150 g m$^{-2}$ factor is the mean interception capacity
per unit LAI \citep{schron2017} and $\theta_\mathrm{litter}$ is an
optional litter-water depth.

\paragraph{Snow-water equivalent.}
Snow on the ground effectively acts as liquid water from the
neutron probe's perspective.
Monthly mean SWE (Eq.~\eqref{eq:swe}) is reported alongside
the month in which snow cover first exceeds
$\SI{150}{\milli\metre}$ SWE — a threshold at which the detector
is largely opaque to cosmic-ray neutrons and calibration is
impractical \citep{bogena2017}.

\paragraph{Soil organic carbon hydrogen equivalent.}
Soil organic matter (SOM) contains hydrogen at a concentration
similar to liquid water.
The volumetric hydrogen equivalent is approximated as
\citep{bogena2017}:
\begin{equation}
  \theta_\mathrm{SOC} \approx
    0.55\,\rho_b\,\omega_\mathrm{SOM}
  \quad[\si{\cubic\metre\per\cubic\metre}],
  \label{eq:thetaSOC}
\end{equation}
where $\omega_\mathrm{SOM}$ is the SOM mass fraction [g/g]
(SOM $\approx 1.724\times$ SOC by convention).
This term is reported alongside $\lw$ as a supplementary offset
that does not change with time.

\subsection{Optimal Soil Sampling Plan}
\label{sec:methods:sampling}

A CRNS calibration requires in-situ gravimetric soil moisture
measurements distributed across the footprint in proportion to
the neutron sensitivity kernel $W(r)$.
The pipeline designs a plan of $n_\mathrm{tot} = 26$ sampling
locations arranged in four concentric rings at radii
$r_k = \{0.25,\, 0.50,\, 0.75,\, 1.00\} \times \rzz$.

\paragraph{Ring weight.}
The contribution of ring $k$ to the footprint integral is
\begin{equation}
  \Phi_k = \int_{r_{k-1}}^{r_k} W(r)\,2\pi r\,\mathrm{d}r,
  \label{eq:ring_weight}
\end{equation}
evaluated numerically by the trapezoidal rule.
Points per ring $n_k \in \{4, 6, 8, 8\}$ are allocated so that
$n_k / \sum n_j \approx \Phi_k / \sum \Phi_j$.

\paragraph{Sample depth.}
At each location, soil is sampled at three depth intervals
(0–5, 5–15, 15–30 cm) to resolve the vertical weighting of
neutron penetration.

\paragraph{Footprint dynamic range.}
As supplementary information, the report includes a plot of
$\rzz(\thetav)$ computed over the full moisture range
$[0.01,\, 0.50]$ \si{\cubic\metre\per\cubic\metre}.
This shows how the footprint radius, and hence the spatial
representativeness of the measurements, changes between dry and
wet conditions at the site.

\subsection{Geological Characterisation}

The Macrostrat API v2 is queried with the site coordinates at three
resolution scales (large, medium, small) in descending order of
coverage.
The first non-empty response is used.
The pipeline extracts:
\begin{itemize}[nosep]
  \item dominant lithology name and Macrostrat rock type
        (carbonate, siliciclastic, igneous, \ldots);
  \item estimated influence on $\lw$ (low, medium, high), based on
        the known hydrogen content of the rock type;
  \item an indicative radiogenicity level (for U/Th-based radiation
        background assessment);
  \item geologic age era (Cenozoic, Mesozoic, Palaeozoic, Precambrian).
\end{itemize}
This qualitative summary complements the quantitative SoilGrids
$\lw$ estimate when geological context is relevant (e.g.\
evaporite-rich or serpentinised terrain).
